% ==========================================================
% ENTREGABLE 3 - INTELIGENCIA DE NEGOCIOS
% Pipeline de Datos y Calidad
% Universidad Estatal Peninsula de Santa Elena
% ==========================================================

\documentclass[12pt,a4paper]{article}

% ==========================================================
% PAQUETES
% ==========================================================
\usepackage[utf8]{inputenc}
\usepackage[T1]{fontenc}
\usepackage[spanish,es-tabla]{babel}
\usepackage[top=2.5cm,bottom=2.5cm,left=2.5cm,right=2.5cm]{geometry}
\usepackage{graphicx}
\usepackage[table]{xcolor}
\usepackage{titlesec}
\usepackage{enumitem}
\usepackage{tabularx}
\usepackage{booktabs}
\usepackage{longtable}
\usepackage{array}
\usepackage{float}
\usepackage{hyperref}
\usepackage{parskip}
\usepackage{pdflscape}

% ==========================================================
% COLORES Y ESTILO
% ==========================================================
\definecolor{primarygray}{gray}{0.35}
\definecolor{textblack}{gray}{0.0}
\definecolor{lightgray}{gray}{0.92}

\titleformat{\section}
{\color{primarygray}\Large\bfseries}
{\thesection.}
{0.5em}
{}

\titleformat{\subsection}
{\color{primarygray}\large\bfseries}
{\thesubsection.}
{0.5em}
{}

\hypersetup{
    colorlinks=true,
    linkcolor=textblack,
    urlcolor=primarygray,
    citecolor=primarygray
}

\renewcommand{\arraystretch}{1.25}
\setlength{\tabcolsep}{4pt}

\newcolumntype{P}[1]{>{\raggedright\arraybackslash}p{#1}}
\newcommand{\StagingTabla}{\texttt{staging.stg\_actividad\_agente\_ia}}
\newcommand{\pathcode}[1]{\begingroup\footnotesize\ttfamily\nolinkurl{#1}\endgroup}

\begin{document}

% ==========================================================
% PORTADA
% ==========================================================
\begin{titlepage}
\centering
\vspace*{0.2cm}

\IfFileExists{upse.jpg}{
    \includegraphics[width=3.8cm]{upse.jpg}\\[0.25cm]
}{
    {\Large \textbf{UPSE}}\\[0.25cm]
}

{\large \textbf{\textcolor{primarygray}{UNIVERSIDAD ESTATAL PEN\'INSULA DE SANTA ELENA}}}\\[0.2cm]
{\normalsize \textbf{Ingenier\'ia de Software}}\\[0.05cm]
{\normalsize Sexto Semestre}\\[0.65cm]

{\normalsize \textbf{VI Inteligencia de Negocios}}\\[0.45cm]
{\Large \textbf{\textcolor{primarygray}{ENTREGABLE 3}}}\\[0.2cm]
{\large Pipeline de Datos y Calidad}\\[0.2cm]
{\normalsize Informe t\'ecnico-acad\'emico de implementaci\'on ETL}\\[0.7cm]

{\normalsize \textbf{T\'itulo del Proyecto}}\\[0.25cm]

\begin{minipage}{0.9\textwidth}
\centering
\normalsize
Plataforma de Inteligencia de Negocios para el an\'alisis de tendencias,
adopci\'on y evoluci\'on de agentes de Inteligencia Artificial en el
dominio del desarrollo de software mediante integraci\'on de datos
heterog\'eneos durante el per\'iodo 2023--2026
\end{minipage}

\vspace{0.45cm}

{\normalsize \textbf{Repositorio del Proyecto}}\\[0.15cm]
{\footnotesize \url{https://github.com/DevTorres404/ai-agents-adoption-observatory.git}}

\vspace{0.55cm}

{\textbf{Autores}}\\[0.2cm]
Damian Jes\'us Torres Cadena\\
Melanie Adriana Tomal\'a Merejildo\\
Jean Pierre Cede\~no Andrade
\vspace{0.45cm}
{\textbf{Docente:}}\\[0.2cm]
Ing. Anthony Pachay
\vspace{0.45cm}

{\textbf{Per\'iodo Acad\'emico}}\\
2026--1

\end{titlepage}

\tableofcontents
\newpage

% ==========================================================
\section{Prop\'osito del Entregable}

El Entregable 3 constituye la evidencia de implementaci\'on del proyecto de Inteligencia de Negocios. A diferencia del Entregable 2, orientado al dise\~no arquitect\'onico y multidimensional, este informe documenta la ejecuci\'on del pipeline ETL, la persistencia de datos en la zona Raw/Bronze, la consolidaci\'on en la zona Staging/Silver y la aplicaci\'on de controles de calidad verificables antes de avanzar hacia el Data Warehouse.

El alcance del entregable comprende la extracci\'on desde fuentes web, API p\'ublica y dataset estructurado; la conservaci\'on de los registros originales; la normalizaci\'on hacia una estructura tabular anal\'itica; y la generaci\'on de evidencias CSV y bit\'acoras persistentes que permiten auditar la trazabilidad del proceso.

% ==========================================================
\section{Arquitectura ETL Implementada}

La soluci\'on implementa una arquitectura tipo Medallion con dos capas operativas principales dentro del alcance de este entregable:

\begin{table}[H]
\centering
\footnotesize
\rowcolors{2}{gray!8}{white}
\begin{tabular}{|P{3cm}|P{5.1cm}|P{5.5cm}|}
\hline
\rowcolor{gray!25}
\textbf{Capa} & \textbf{Implementaci\'on} & \textbf{Funci\'on dentro del E3} \\ \hline
Raw / Bronze & Archivos JSON por fuente y tablas \pathcode{raw.raw_files}, \pathcode{raw.raw_records} en PostgreSQL & Conserva la informaci\'on original de cada fuente sin transformarla indebidamente. \\ \hline
Staging / Silver & Tabla \StagingTabla{} reconstruida mediante scripts Python & Integra, limpia, normaliza y homologa los registros para dejarlos listos para el Data Warehouse. \\ \hline
Quality & M\'odulos \pathcode{src/quality} y evidencias en \pathcode{docs/evidencias} & Verifica completitud, duplicados, nulos, casting, homologaci\'on, trazabilidad y errores de ejecuci\'on. \\ \hline
\end{tabular}
\caption{Capas t\'ecnicas implementadas en el Entregable 3}
\end{table}

El script principal de orquestaci\'on es \pathcode{src/scripts/run_pipeline.py}. La capa Staging unificada se reconstruye desde \pathcode{src/staging/stg_build_unified.py}, lo que permite reproducir el flujo a partir de los registros conservados en Raw.

% ==========================================================
\section{Fuentes de Datos Implementadas}

El proyecto integra siete fuentes principales para observar actividad t\'ecnica, producci\'on especializada, popularidad, innovaci\'on, percepci\'on comunitaria y adopci\'on acad\'emica sobre agentes de Inteligencia Artificial aplicados al desarrollo de software. La fuente propia UPSE se incorpora desde \pathcode{data/encuesta/encuesta.json}, correspondiente a respuestas reales de Google Forms.

\scriptsize
\begin{longtable}{|P{2.4cm}|P{2.2cm}|P{2.4cm}|P{6.1cm}|}
\hline
\rowcolor{gray!25}
\textbf{Fuente} & \textbf{M\'etodo} & \textbf{Rol Anal\'itico} & \textbf{Evidencia t\'ecnica esperada} \\ \hline
\endfirsthead
\hline
\rowcolor{gray!25}
\textbf{Fuente} & \textbf{M\'etodo} & \textbf{Rol Anal\'itico} & \textbf{Evidencia t\'ecnica esperada} \\ \hline
\endhead
Reddit & Web scraping & Comunidad & Publicaciones y menciones relacionadas con agentes de IA en comunidades de desarrolladores. Raw en \pathcode{data/raw/scraping/reddit/}. \\ \hline
DEV Community & Web scraping / API p\'ublica & Producci\'on t\'ecnica & Art\'iculos, tutoriales y experiencias t\'ecnicas de desarrolladores. Raw en \pathcode{data/raw/scraping/devto/}. \\ \hline
Hacker News & Web scraping & Innovacion tecnologica & Discusiones y tendencias emergentes. Raw en \pathcode{data/raw/scraping/hackernews/}. \\ \hline
Google Trends & Web scraping / pytrends & Popularidad & Interes relativo de busqueda por agente y periodo. Raw en \pathcode{data/raw/scraping/google_trends/}. \\ \hline
GitHub & API p\'ublica REST & Actividad t\'ecnica & Repositorios, estrellas, forks, issues y actividad asociada a agentes de IA. Raw en \pathcode{data/raw/api/github/}. \\ \hline
AIDev Dataset: AI Coding & Archivo estructurado Parquet & Catalogo maestro & Dataset estructurado convertido a Raw JSON analitico. Raw en \pathcode{data/raw/archivos/catalogo/}. \\ \hline
Fuente propia UPSE & Google Forms / JSON & Adopci\'on acad\'emica & Respuestas reales normalizadas desde \pathcode{data/encuesta/encuesta.json}. Raw en \pathcode{data/raw/fuente_propia/fuente_propia/}. \\ \hline
\caption{Fuentes de datos utilizadas en el Entregable 3}
\end{longtable}
\normalsize

\subsection{Par\'ametros Din\'amicos de Extracci\'on}

\begin{table}[H]
\centering
\scriptsize
\rowcolors{2}{gray!8}{white}
\begin{tabular}{|P{2.7cm}|P{4.5cm}|P{6.2cm}|}
\hline
\rowcolor{gray!25}
\textbf{Fuente} & \textbf{Script} & \textbf{Parametros dinamicos o configurables} \\ \hline
GitHub REST API & \pathcode{src/extractors/github_api.py} & \pathcode{queries}, \pathcode{pages}, \pathcode{per_page}, token opcional y rango temporal desde 2023. \\ \hline
Hacker News & \pathcode{src/extractors/hackernews_scraper.py} & URL objetivo, l\'imite de registros y metadatos de ejecuci\'on. \\ \hline
Dev.to & \pathcode{src/extractors/devto_scraper.py} & Busquedas por terminos de agentes, limite de articulos, filtro de relevancia y rango temporal desde 2023. \\ \hline
Reddit & \pathcode{src/extractors/reddit_scraper.py} & B\'usquedas t\'ecnicas, l\'imite de posts, selectores Playwright y metadata temporal. \\ \hline
Google Trends & \pathcode{src/extractors/trends_scraper.py} & Lista de keywords, region, timeframe \pathcode{2023-01-01 2026-12-31} y manejo de rate limit. \\ \hline
AIDev Dataset & \pathcode{src/extractors/aidedev_catalog.py} & Rutas Parquet, descarga controlada, agentes homologados y filtro temporal 2023--2026. \\ \hline
Fuente propia UPSE & \pathcode{src/extractors/google_forms_survey.py} & Archivo \pathcode{data/encuesta/encuesta.json}, normalizaci\'on de encabezados y rango temporal 2023--2026. \\ \hline
\end{tabular}
\caption{Matriz de par\'ametros de extracci\'on}
\end{table}

% ==========================================================
\section{Zona Raw / Bronze}

La zona Raw conserva los registros originales en archivos por fuente y en PostgreSQL usando campos JSONB. La convenci\'on cronol\'ogica definida para el entregable es:

\begin{center}
\texttt{fuente\_YYYY-MM-DD.ext}
\end{center}

Ejemplos esperados de archivos Raw:

\begin{verbatim}
data/raw/api/github/github_2026-06-28.json
data/raw/scraping/devto/devto_2026-06-28.json
data/raw/scraping/hackernews/hackernews_2026-06-28.json
data/raw/scraping/reddit/reddit_2026-06-28.json
data/raw/scraping/google_trends/google_trends_2026-06-28.json
data/raw/archivos/catalogo/aidedev_2026-06-28.json
data/raw/fuente_propia/fuente_propia/fuente_propia_2026-07-03.json
\end{verbatim}

El inventario completo de Raw no debe calcularse solo con archivos de la \'ultima corrida, sino desde PostgreSQL. La evidencia correspondiente se exporta en:

\begin{itemize}[leftmargin=1.2cm]
    \item \pathcode{docs/evidencias/inventario_raw.csv}
    \item \pathcode{docs/evidencias/inventario_raw_completo.csv}
\end{itemize}

% ==========================================================
\section{Zona Staging / Silver}

La tabla de Staging consolidada es \StagingTabla{}. Su objetivo es convertir datos heterog\'eneos en una estructura tabular consistente para el modelo dimensional del Entregable 4.

\begin{table}[H]
\centering
\footnotesize
\rowcolors{2}{gray!8}{white}
\begin{tabular}{|P{4cm}|P{9.2cm}|}
\hline
\rowcolor{gray!25}
\textbf{Aspecto} & \textbf{Contrato documentado} \\ \hline
Tabla f\'isica & 23 columnas en PostgreSQL. \\ \hline
Columnas anal\'iticas & 21 columnas orientadas al an\'alisis. \\ \hline
Columnas t\'ecnicas & \pathcode{id} y \pathcode{fecha_carga}. \\ \hline
Evidencia & \pathcode{docs/evidencias/staging_contract_columns.csv}. \\ \hline
\end{tabular}
\caption{Contrato actual de Staging}
\end{table}

Las transformaciones principales implementadas son:

\begin{itemize}[leftmargin=1.2cm]
    \item Aplanamiento de estructuras JSON para obtener columnas anal\'iticas uniformes.
    \item Normalizaci\'on de fechas ISO y filtrado del per\'iodo 2023--2026.
    \item Clasificaci\'on de agentes mediante reglas de expresiones regulares.
    \item Categorizaci\'on estrat\'egica seg\'un rol anal\'itico de la fuente y tipo de se\~nal.
    \item Deduplicaci\'on por clave compuesta.
\end{itemize}

La regla de deduplicaci\'on aplicada es:

\begin{center}
\texttt{fuente + plataforma + id\_origen\_registro + nombre\_agente + fecha\_evento}
\end{center}

Esta regla evita contar varias veces el mismo evento cuando aparece duplicado dentro de una fuente o por recargas controladas.

% ==========================================================
\section{Framework de Calidad de Datos}

El framework de calidad implementa los siete controles obligatorios del Entregable 3. Cada control genera evidencia verificable mediante CSV, logs o consultas sobre PostgreSQL.

\scriptsize
\begin{longtable}{|P{3.7cm}|P{4.8cm}|P{4.8cm}|}
\hline
\rowcolor{gray!25}
\textbf{Control} & \textbf{Implementaci\'on} & \textbf{Evidencia} \\ \hline
\endfirsthead
\hline
\rowcolor{gray!25}
\textbf{Control} & \textbf{Implementaci\'on} & \textbf{Evidencia} \\ \hline
\endhead
Completitud Raw/Staging & Comparaci\'on entre registros Raw acumulados y registros consolidados en Staging. & \pathcode{quality_summary.csv}. \\ \hline
Duplicados reales & Conteo por clave compuesta de deduplicaci\'on. & \pathcode{dedup_report.csv}. \\ \hline
Nulos cr\'iticos & Revisi\'on de columnas obligatorias para an\'alisis. & \pathcode{nulls_matrix.csv}. \\ \hline
Formato y casting & Validaci\'on de fechas, num\'ericos y transformaciones. & \pathcode{casting_report.csv}. \\ \hline
Homologaci\'on entre fuentes & Mapeo de nombres de agentes, plataformas y categor\'ias. & \pathcode{homologation_map.csv}. \\ \hline
Bit\'acora de errores & Registro persistente de fallos HTTP, rate limit y anomal\'ias. & \pathcode{logs/pipeline_run.log}, \pathcode{logs/pipeline_errors.csv}. \\ \hline
M\'etricas finales & Resumen estad\'istico reproducible. & \pathcode{staging_stats.csv}, \pathcode{quality_issue_breakdown.csv}. \\ \hline
\caption{Controles de calidad implementados}
\end{longtable}
\normalsize

% ==========================================================
\section{M\'etricas Verificables}

Las m\'etricas se reportan sobre el universo Raw acumulado en PostgreSQL. En arquitectura Medallion, Raw/Bronze conserva el historial de cargas y Silver/Staging consolida el conjunto anal\'itico mediante deduplicaci\'on controlada.

\begin{table}[H]
\centering
\footnotesize
\rowcolors{2}{gray!8}{white}
\begin{tabular}{|P{7.5cm}|P{3.8cm}|}
\hline
\rowcolor{gray!25}
\textbf{M\'etrica} & \textbf{Valor verificado} \\ \hline
Registros Raw/Bronze acumulados & 358440 \\ \hline
Registros Staging/Silver consolidados & 120846 \\ \hline
Completitud Raw acumulado vs Staging & 33.71\% \\ \hline
Merma controlada por deduplicaci\'on hist\'orica & 66.29\% \\ \hline
Duplicados reales removidos & 237538 \\ \hline
Registros fuente propia Google Forms & 12 \\ \hline
Nulos cr\'iticos & 0 \\ \hline
\end{tabular}
\caption{M\'etricas finales respaldadas por evidencia CSV}
\end{table}

Estas cifras deben contrastarse con los archivos de evidencia generados al finalizar la corrida del pipeline. Si una fuente externa falla por restricciones de red, cambios HTML o rate limit, el evento debe quedar documentado en logs y no debe presentarse como extracci\'on exitosa de la \'ultima corrida.

% ==========================================================
\section{Evidencias Generadas}

Los archivos principales de auditor\'ia t\'ecnica se encuentran en \pathcode{docs/evidencias/}:

\begin{itemize}[leftmargin=1.2cm]
    \item \pathcode{source_execution_evidence.csv}
    \item \pathcode{inventario_raw.csv}
    \item \pathcode{inventario_raw_completo.csv}
    \item \pathcode{staging_stats.csv}
    \item \pathcode{quality_summary.csv}
    \item \pathcode{quality_issue_breakdown.csv}
    \item \pathcode{aidedev_agent_summary.csv}
    \item \pathcode{dedup_report.csv}
    \item \pathcode{nulls_matrix.csv}
    \item \pathcode{casting_report.csv}
    \item \pathcode{homologation_map.csv}
    \item \pathcode{staging_contract_columns.csv}
\end{itemize}

% ==========================================================
\section{Conclusiones}

El Entregable 3 queda documentado como una implementaci\'on reproducible y defendible del pipeline de datos. La arquitectura Raw/Bronze y Staging/Silver est\'a separada, los extractores generan evidencia por fuente, la capa Raw conserva datos originales, Staging ofrece una estructura tabular preparada para el Data Warehouse y el framework de calidad produce archivos auditables.

La fuente propia UPSE queda integrada como fuente acad\'emica real mediante Google Forms. Con la capa Silver consolidada, el proyecto puede avanzar al Entregable 4 tomando \StagingTabla{} como base para construir la Capa Gold del modelo dimensional.

\end{document}
